%================================
% Vocabulary & Language Focus — XB3 (Unit 2)
% 编译方式: xelatex main.tex (需运行两遍以生成目录)
%================================
\documentclass[../main.tex]{subfiles}

\begin{document}

\section{Unit 2}

\vw{gravity}{/ˈɡrævəti/}{n.}

\sense{n.}{重力}
\col{the law of gravity}
\ex{According to the law of gravity, everything that goes up must come down.}
\col{zero gravity 失重}
\ex{Astronauts experience zero gravity when they orbit the earth.}

\sense{n.}{严重性}
\col{the gravity of a situation/offence}
\ex{The judge emphasised the gravity of the offence before passing sentence.}

\sense{n.}{严肃}
\col{with gravity}
\ex{The headmaster addressed the students with great gravity.}

\begin{origin}
借自拉丁语 \textit{gravitatem}（重量、严重性），是 \textit{gravis}（重的、严肃的）的抽象名词。同源词 \textbf{grave}（严重的；严肃的）、\textbf{grief}（悲伤；\textit{gravis} 的罗曼语形式）、\textbf{aggravate}（加重；激怒）、\textbf{gravid}（怀孕的；\textit{gravidus} 充满重量的）。"重力"与"严重性"在拉丁词源中同根——对古罗马人而言，沉重即是严重、严肃、悲伤。
\end{origin}

\vw{absence}{/ˈæbsəns/}{n.}

\sense{n.}{缺乏，不存在}
\col{in the absence of sth}
\ex{In the absence of clear evidence, the court had to acquit the defendant.}
\col{absence of mind 心不在焉}
\ex{His frequent absence of mind worried his colleagues.}

\sense{n.}{缺席}
\col{absence from sth}
\ex{His repeated absence from class led to a formal warning.}
\col{be absent from sth}
\ex{She was absent from work for two weeks with the flu.}

\begin{gram}{the absence of / the presence of：抽象"有无"对照}
英语学术语体高频句式：\textbf{the presence of sth}（……的存在）vs. \textbf{the absence of sth}（……的缺乏）。当 absence 后接抽象名词时，常无冠词：\textit{in the absence of any concrete evidence}。同构：\textit{a complete absence of interest}（毫无兴趣）、\textit{a conspicuous absence}（明显缺席）。
\end{gram}

\begin{origin}
15 世纪借自拉丁语 \textit{absentia}（不在场），是 \textit{abesse}（不在）的派生形式：\textit{ab-}（远离）+ \textit{esse}（存在）。同源词 \textbf{absent}（缺席的；不在的）、\textbf{absently}（心不在焉地）、\textbf{absent-minded}（心不在焉的）、\textbf{essence}（本质；\textit{esse} $\to$存在之物）、\textbf{entity}（实体）。absence 与 present（出席）相对，与 absolute（绝对的）部分同源。
\end{origin}

\vw{float}{/fləʊt/}{vi. \& vt.}

\sense{vi.}{飘动，漂流}
\col{float in/on the water/air}
\ex{Empty bottles floated on the surface of the lake.}
\col{float about/around}
\ex{Papers were floating about the room after the burglary.}

\sense{vi.}{浮}
\col{float on sth}
\ex{Wood floats on water, but iron does not.}
\col{floating dock/floating rate 浮船坞/浮动汇率}
\ex{The company operates with a floating exchange rate.}

\sense{vt.}{使浮动，使漂流}
\col{float a plan/idea/scheme}
\ex{The minister floated the idea of a new tax on luxury goods.}
\col{float a company 上市（让公司股票开始交易）}
\ex{The company was floated on the stock exchange in 2010.}

\begin{gram}{float a loan / float an idea：抽象"投放"用法}
\textbf{float a loan}（发行贷款）、\textbf{float a currency}（让货币汇率自由浮动）、\textbf{float a rumour}（散布谣言）。float 抽象用法皆源自物理"漂浮"——贷款/货币/谣言"在市场上漂浮"，等待接收者"接住"。
\end{gram}

\begin{origin}
古英语 \textit{flotian}（漂浮、游泳），与 \textit{fleotan}（流、漂浮）同源，原始日耳曼语 *\textit{flot-}。同源词 \textbf{fleet}（船队；流过；\textit{fleotan} 的派生）、\textbf{flotsam}（漂浮残骸；\textit{float} + 古英语 \textit{sam} 集合）、\textbf{flood}（洪水；\textit{fleotan} $\to$ 流）。float 与 flow、flush、fluster 同根——皆"流动"族。
\end{origin}

\vw{cupboard}{/ˈkʌbəd/}{n.}

\sense{n.}{壁橱，橱柜，衣柜}
\col{a kitchen/wall/built-in cupboard}
\ex{She opened the kitchen cupboard and took out a jar of coffee.}
\col{a skeleton in the cupboard 家丑（英式）}
\ex{Every family has a skeleton in the cupboard that no one talks about.}

\begin{gram}{a skeleton in the cupboard / closet：英美文化差异}
英式英语惯说 \textbf{skeleton in the cupboard}（橱柜里的骷髅），美式英语惯说 \textbf{skeleton in the closet}（壁橱里的骷髅）——指不可告人的家丑。同义、异国：美式英语"closet"一词既指橱柜（口语）也指衣橱（书面）。
\end{gram}

\begin{origin}
15 世纪由 \textit{cup}（杯；木盆）+ \textit{board}（板）合成，字面义"放杯子的板"。cup 借自拉丁语 \textit{cuppa}（杯），board 来自古英语 \textit{bord}（板）。同族还有 \textbf{cupboard love}（因利益而表的爱，"为得到糖果而讨好"——小孩对父母的"橱柜之爱"）。
\end{origin}

\vw{ceiling}{/ˈsiːlɪŋ/}{n.}

\sense{n.}{天花板}
\col{a high/low ceiling}
\ex{The living room has a high ceiling and tall windows.}
\col{ceiling fan/light 吊扇/吸顶灯}
\ex{The ceiling fan slowly turned in the afternoon heat.}

\sense{n.}{上限}
\col{a ceiling on sth}
\ex{The government has imposed a ceiling on immigration numbers.}
\col{glass ceiling 玻璃天花板（无形的职业晋升障碍）}
\ex{Women still face a glass ceiling in many senior corporate positions.}

\begin{lex}{ceiling / roof / floor / ground 的"上下限"对照}
\textbf{ceiling}（上限；天花板）；\textbf{roof}（屋顶；不可超越的"顶层"）；\textbf{floor}（下限；地板）；\textbf{ground}（地面；最低基线）。在抽象语境中：\textit{price ceiling/floor}（价格上限/下限）、\textit{wage ceiling}（工资上限）、\textit{noise floor}（噪声底线）。
\end{lex}

\begin{origin}
14 世纪由 \textit{ceil}（装顶板；来自古法语 \textit{ciel} 天花板、天空）+ \textit{-ing} 派生。\textit{ciel} 来自拉丁语 \textit{caelum}（天空）。同源词 \textbf{ceiling} 的"上限"义源自物理天花板——不可再高之限。
\end{origin}

\vw{bother}{/ˈbɒðə(r)/}{vt. \& vi. \& n.}

\sense{vt.}{使烦恼，使担心}
\col{bother sb}
\ex{I'm sorry to bother you, but could you help me with this?}
\col{bother sb with/about sth}
\ex{Don't bother your mother with such trivial questions.}

\sense{vt.}{花费时间精力（做某事）}
\col{bother to do sth}
\ex{He didn't even bother to reply to my email.}
\col{bother doing sth}
\ex{Why bother asking if you're not going to listen to the answer?}

\sense{vt.}{打扰}
\col{bother sb}
\ex{I'm busy — don't bother me now.}

\sense{n.}{麻烦，困难}
\col{a bit of bother}
\ex{It's no bother at all — I'm happy to help.}
\col{have bother (doing sth)}
\ex{We had a bit of bother finding the place.}

\begin{gram}{bother / disturb / trouble / annoy 的"烦扰"强度}
\textbf{bother} 轻度烦恼（小事打扰：\textit{Sorry to bother you}）；\textbf{disturb} 强调打断（安静/睡眠：\textit{Please don't disturb the baby}）；\textbf{trouble} 给别人添麻烦（正式：\textit{May I trouble you for the salt?}）；\textbf{annoy} 持续恼怒（\textit{His humming really annoys me}）。bother 还可不接宾语单独使用：\textit{Don't bother.}（不必麻烦了）。
\end{gram}

\begin{origin}
17 世纪方言爱尔兰英语，可能与爱尔兰语 \textit{bodhar}（聋的、烦躁的）相关。同源词 \textbf{bothersome}（烦人的）、\textbf{bothered}（烦恼的；英式口语中表"愿意的"：\textit{I'm not bothered whether we go or not} 我无所谓）。bother 是 18 世纪后英语中替代 trouble 的口语化选择。
\end{origin}

\vw{occupy}{/ˈɒkjupai/}{vt.}

\sense{vt.}{占用，使用}
\col{occupy sth}
\ex{The new sofa occupies most of the living room.}
\col{occupy time/space}
{Reading occupies most of my spare time.}

\sense{vt.}{居住}
\col{occupy a house/flat}
\ex{The family has occupied the flat for over twenty years.}

\sense{vt.}{占领，侵占}
\col{occupy a country/city}
\ex{The army occupied the capital within three days.}
\col{an occupied zone/territory}
\ex{Life in the occupied zone was harsh and unpredictable.}

\sense{vt.}{使忙于}
\col{occupy oneself with/in doing sth}
\ex{During the long winter evenings, she occupied herself with painting and reading.}
\col{be occupied with sth}
\ex{The staff are fully occupied with the new project.}

\begin{gram}{occupied / vacant / engaged 的"占/空"细分}
\textbf{occupied}（已被占用：\textit{the seat is occupied}）；\textbf{vacant}（空着：\textit{a vacant room}）；\textbf{engaged}（正被使用，且通常指短时间内不可用：英式 \textit{engaged} = 美式 \textit{busy}，多用于电话/洗手间）。英式英语"toilet is engaged"即美式"toilet is busy"。
\end{gram}

\begin{origin}
16 世纪借自拉丁语 \textit{occupare}（占据、占有）：\textit{ob-}（加强前缀，朝向）+ \textit{capere}（抓、拿）。同源词 \textbf{occupation}（职业；占领）、\textbf{occupant}（居住者）、\textbf{preoccupied}（心事重重的；\textit{prae-} 在前 + occupare $\to$ 心思先被占）、\textbf{capture}（俘获）、\textbf{capable}（有能力的）、\textbf{capacity}（容量）、\textbf{captive}（俘虏）。
\end{origin}

\vw{mission}{/ˈmɪʃn/}{n.}

\sense{n.}{任务}
\col{on a mission}
\ex{The soldiers were sent on a dangerous mission behind enemy lines.}
\col{carry out/complete a mission}
\ex{The team successfully carried out the mission in less than 24 hours.}

\sense{n.}{使命，天职}
\col{a mission in life}
\ex{She regarded teaching as her mission in life.}
\col{a man/woman with a mission 有使命感的人}
\ex{He's a man with a mission — once he sets his mind on something, nothing can stop him.}

\sense{n.}{军事行动；太空飞行任务}
\col{a military/rescue/combat mission}
\ex{The rescue mission saved over a hundred lives.}
\col{space mission}
\ex{The first manned mission to Mars is scheduled for 2030.}

\begin{lex}{mission / task / assignment / duty 的"任务"细分}
\textbf{mission} 强调重要、长期、带有使命感（\textit{a peace mission}）；\textbf{task} 通用，可大可小（\textit{a difficult task}）；\textbf{assignment} 多指工作/学业布置（\textit{a homework assignment}）；\textbf{duty} 强调责任（\textit{It's my duty to report this}）。mission 比 task 更崇高，比 duty 更具体。
\end{lex}

\begin{origin}
16 世纪借自拉丁语 \textit{missio}（释放、派出），源于 \textit{mittere}（发送、释放）。同源词 \textbf{missionary}（传教士）、\textbf{missile}（导弹；"被发射之物"）、\textbf{mission statement}（使命宣言）、\textbf{admit}（承认；准入）、\textbf{emit}（发射；发出）、\textbf{submit}（提交；屈服）、\textbf{transmit}（传输）、\textbf{promise}（承诺；\textit{pro-} 前 + \textit{mittere} $\to$预先放出之言）。"任务"原义"被派出去"。
\end{origin}

\vw{evaluate}{/ɪˈvæljueɪt/}{vt.}

\sense{vt.}{评估，估计，评价}
\col{evaluate sth}
\ex{Teachers evaluate students' progress through a combination of tests and projects.}
\col{evaluate the situation}
\ex{Before deciding, we need to evaluate the situation more carefully.}

\begin{gram}{evaluate / assess / estimate / appraise 的"估"细分}
\textbf{evaluate} 强调权衡多因素后给出价值判断（\textit{evaluate the performance}）；\textbf{assess} 强调定量/定性估价（\textit{assess the damage}）；\textbf{estimate} 强调根据有限信息粗略估计（\textit{estimate the cost}）；\textbf{appraise} 强调专业/官方评估（\textit{appraise the property}）。evaluate 偏综合判断，estimate 偏数据估算。
\end{gram}

\begin{origin}
19 世纪借自法语 \textit{évaluer}，法语回追至古法语 \textit{value}（价值），拉丁语 \textit{valere}（强、有价值）。同源词 \textbf{value}（价值）、\textbf{valuation}（估价）、\textbf{valuable}（贵重的）、\textbf{invaluable}（无价的）、\textbf{valid}（有效的）、\textbf{prevail}（盛行；\textit{prae-} 极 + \textit{valere}）、\textbf{equivalent}（等价的）。
\end{origin}

\vw{creature}{/ˈkriːtʃə(r)/}{n.}

\sense{n.}{动物，生物}
\col{creature of ...}
\ex{The deep ocean is home to many strange creatures that have never been seen by human eyes.}
\col{a living creature}
\ex{All living creatures deserve our respect and protection.}

\sense{n.}{人（带有感情色彩）}
\col{a poor/strange/lazy creature}
\ex{The poor creature had been wandering in the snow for hours.}
\col{creature of habit 习惯的奴隶}
\ex{My father is a creature of habit — he eats the same breakfast every day.}

\begin{gram}{creature comfort：物质享乐的"造物慰藉"}
\textbf{creature comforts}（物质享受；使身体舒适之物：\textit{a hotel with all the creature comforts} 配备一切舒适设施的酒店）。comfort 加上 creature 是修辞性的拟人化——动物（人）所需之慰藉。
\end{gram}

\begin{origin}
13 世纪借自古法语 \textit{creature}，拉丁语 \textit{creatura}（被创造之物），是 \textit{creare}（创造）的派生形式。同源词 \textbf{create}（创造）、\textbf{creation}（创造；作品）、\textbf{creative}（创造性的）、\textbf{creator}（创造者）、\textbf{recreate}（再现；娱乐）、\textbf{procreate}（生育）、\textbf{concrete}（具体的；\textit{con-} 一起 + \textit{crescere} 长$\to$长在一起）。"造物"在神学语境特指"上帝所造之物"——creature comfort 词源上即"上帝创造之物的慰藉"。
\end{origin}

\vp{cast new light on}

\sense{短语}{使进一步了解……}
\col{cast/shed new light on sth}
\ex{Recent discoveries have cast new light on the causes of the disease.}
\ex{Her autobiography casts an interesting new light on the political events of the 1960s.}

\begin{gram}{cast/shed/throw light on：光之三态动词}
\textbf{cast light on}（投光$\to$阐明，最正式）、\textbf{shed light on}（散光$\to$揭示，最常用）、\textbf{throw light on}（抛光$\to$说明，最口语）。三短语中 light 后常接 on/upon + 待阐明之事。同族：\textit{in the light of}（鉴于；考虑到）、\textit{come to light}（被发现）、\textit{bring to light}（揭露）。
\end{gram}

\vw{maintenance}{/ˈmeɪntənəns/}{n.}

\sense{n.}{维护，保养}
\col{car/road maintenance}
\ex{The bridge has been closed for routine maintenance.}
\col{maintenance costs}
\ex{The new design significantly reduces long-term maintenance costs.}

\sense{n.}{维持，保持}
\col{the maintenance of sth}
\ex{The maintenance of law and order is the primary duty of the police.}
\col{maintenance of good relations}
\ex{The two countries attach great importance to the maintenance of good diplomatic relations.}

\begin{origin}
15 世纪借自古法语 \textit{maintenance}，源于 \textit{maintenir}（维持、支持），拉丁语 \textit{manu tenere}（用手拿住）：\textit{manus}（手）+ \textit{tenere}（持）。同源词 \textbf{maintain}（维持）、\textbf{maintainable}（可维持的）、\textbf{abstain}（戒绝；\textit{ab-} 离 + \textit{tenere}）。maintenance 字面"手拿住"——握住不放即是维持。
\end{origin}

\vw{circumstance}{/ˈsɜːkəmstəns/}{n.}

\sense{n.}{条件，状况}
\col{under/in the circumstances}
\ex{Under the circumstances, I think we did the right thing.}
\col{under no circumstances 绝不}
\ex{Under no circumstances should you open this door.}

\sense{n.}{(circumstances) 境况，(尤指) 经济状况}
\col{family/financial/personal circumstances}
\ex{Students from less well-off family circumstances receive a grant from the university.}
\col{circumstances permit/allow}
\ex{If circumstances permit, I'll join you for the weekend.}

\begin{gram}{under no circumstances 后接倒装句}
\textbf{Under no circumstances} 后接助动词提前的倒装句：\textit{Under no circumstances should you lend money to strangers.}（\ding{55} \textit{Under no circumstances you should lend...}）。同类倒装：\textit{Never have I seen such a thing.}、\textit{Not only did he apologise, but he also paid for the damage.}。否定意义副词提前至句首是英语"否定倒装"经典句式。
\end{gram}

\begin{origin}
13 世纪借自古法语 \textit{circonstance}，拉丁语 \textit{circumstantia}（周围的情形），是 \textit{circumstare}（站立于四周）的派生形式：\textit{circum-}（周围）+ \textit{stare}（站）。同源词 \textbf{circumstantial}（视情况而定的；详细的）、\textbf{circumvent}（绕行；规避；\textit{circum-} 周围 + \textit{venire} $\to$绕到……之后）、\textbf{stance}（姿态；立场）、\textbf{constant}（恒定的；\textit{con-} 一起 + stare $\to$ 一起站）。"境况"字面是"站在四周之物"。
\end{origin}

\vp{in the circumstances}

\sense{短语}{在这种情况下，既然如此}
\ex{In the circumstances, I think we should postpone the meeting.}
\ex{Given the bad weather, in the circumstances it is best to stay indoors.}

\begin{gram}{in / under the circumstances / given the circumstances：三态同义}
三者皆表"鉴于这种情况"：\textbf{in the circumstances}（英式正式）、\textbf{under the circumstances}（英美通用，最常见）、\textbf{given the circumstances}（美式偏好）。细微差别：in 强调"处于这种环境之中"，under 强调"在这种条件之下"，given 强调"已经知道这些事实"。
\end{gram}

\vw{exposure}{/ɪkˈspəʊʒə(r)/}{n.}

\sense{n.}{面临，遭受}
\col{exposure to sth}
\ex{Prolonged exposure to loud noise can cause permanent hearing damage.}
\col{exposure to radiation}
\ex{The workers were monitored for exposure to radiation.}

\sense{n.}{揭露}
\col{the exposure of sth}
\ex{The exposure of the corruption scandal led to the minister's resignation.}

\sense{n.}{报道}
\col{media exposure}
\ex{The film received massive media exposure during the awards season.}
\col{public exposure}
\ex{She won the award for her outstanding public exposure of the human rights issue.}

\begin{lex}{exposure / disclosure / revelation 的"揭露"细分}
\textbf{exposure} 强调"暴露"——使原本隐藏的东西显形（\textit{exposure to the cold} 受寒；\textit{exposure of the truth} 真相暴露）；\textbf{disclosure} 强调"披露"——主动公开秘密（\textit{the disclosure of classified information} 机密公开）；\textbf{revelation} 强调"惊人的揭露"——常带戏剧性（\textit{his revelations shocked the country}）。
\end{lex}

\begin{origin}
17 世纪由 \textit{expose}（暴露）派生：古法语 \textit{exposer}（暴露），拉丁语 \textit{exponere}（摆出、暴露）：\textit{ex-}（出）+ \textit{ponere}（放置）。同源词 \textbf{expose}（暴露）、\textbf{exposé}（揭露材料；法语形式）、\textbf{compose}（组成；作曲；\textit{com-} 一起 + ponere）、\textbf{dispose}（处理；安排）、\textbf{oppose}（反对；\textit{op-} 对 + ponere）、\textbf{propose}（提议；求婚；\textit{pro-} 前 + ponere $\to$放在前面）、\textbf{impose}（强加；征税）。
\end{origin}

\vw{visible}{/ˈvɪzəbl/}{adj.}

\sense{adj.}{看得见的}
\col{visible to sb}
\ex{The stars are clearly visible to the naked eye on a clear night.}
\col{clearly/barely visible}
\ex{The tiny island was barely visible on the horizon.}

\sense{adj.}{明显的}
\col{a visible improvement/difference/effort}
\ex{There has been a visible improvement in her health since she started the new treatment.}

\begin{gram}{visible / visual / visionary / vision 的派生分叉}
\textbf{visible}（看得见的；形容词）、\textbf{visual}（视觉的；多用于名词 \textit{visual aids} 视觉教具）、\textbf{vision}（视力；视野；远见）、\textbf{visionary}（有远见的；空想的）、\textbf{invisible}（看不见的）。注意 vision 与 visible 同源但词义已分——vision 偏"看见的能力"（精神/物理），visible 偏"可以被看见"。
\end{gram}

\begin{origin}
14 世纪借自拉丁语 \textit{visibilis}（看得见的），源于 \textit{videre}（看）。同源词 \textbf{vision}（视力；视野）、\textbf{visual}（视觉的）、\textbf{visage}（面容；\textit{vis-} 看的变体）、\textbf{visit}（参观；"去看"）、\textbf{revise}（修订；复习；\textit{re-} 再 + \textit{visere} 看）、\textbf{supervise}（监督；\textit{super-} 上 + videre）、\textbf{evident}（明显的；\textit{e-} 出 + \textit{videre} $\to$ 看得出来的）、\textbf{provide}（提供；\textit{pro-} 前 + \textit{videre} $\to$ 预先看到$\to$预备）。
\end{origin}

\vp{in view of}

\sense{短语}{鉴于，考虑到}
\ex{In view of the bad weather, the match has been postponed.}
\ex{In view of her long service to the company, the directors gave her a special award.}

\begin{lex}{in view of / in case of / owing to / on account of 的"鉴于"细分}
\textbf{in view of}（鉴于；考虑到，正面论证：\textit{In view of his age, we excused him}）；\textbf{in case of}（万一；防备：\textit{In case of fire, break the glass}）；\textbf{owing to}（由于；归因：\textit{Owing to the rain, the match was cancelled}）；\textbf{on account of}（因为；归因：\textit{on account of his illness}）。in view of 表"考虑到"以做合理判断，in case of 表"以防万一"。
\end{lex}

\vw{luxury}{/ˈlʌkʃəri/}{n.}

\sense{n.}{不常有的乐趣（或享受）}
\col{a luxury of sth}
\ex{A long, hot bath after a hard day's work is a luxury I really appreciate.}
\col{luxuries of modern life}
\ex{Electricity, clean water, and central heating are luxuries that many people still take for granted.}

\sense{n.}{奢侈品}
\col{a luxury car/hotel/item}
\ex{She bought herself a little luxury — a silk scarf from Paris.}
\col{luxury goods}
\ex{The shop specialises in luxury goods such as handbags, watches and fine wine.}

\begin{gram}{luxury / luxurious / luxuriant 的细微差异}
\textbf{luxury}（n. 奢侈；奢侈品；形容词性：\textit{a luxury hotel}）、\textbf{luxurious}（adj. 奢侈的；豪华的：\textit{a luxurious lifestyle}）、\textbf{luxuriant}（adj. 茂盛的；丰富的：\textit{luxuriant vegetation} 茂盛的植被；不可表"豪华"）。后三者字形相近、词义有别。
\end{gram}

\begin{origin}
14 世纪借自古法语 \textit{luxurie}，拉丁语 \textit{luxuria}（奢侈、过度），源于 \textit{luxus}（过度、奢华）。同源词 \textbf{luxurious}（豪华的）、\textbf{luxuriate}（纵情享受）、\textbf{luxuriant}（茂盛的；同源但分支）、\textbf{luxury tax}（奢侈品税）。luxus 与 \textit{lux}（光）同源——奢侈是"光芒过盛之物"。
\end{origin}

\vw{tube}{/tjuːb/}{n.}

\sense{n.}{管，管子}
\col{a tube of ...}
\ex{She applied a tube of toothpaste to her brush.}
\col{a test tube}
\ex{The scientist poured the liquid from one test tube to another.}

\sense{n.}{软管}
\col{a rubber/plastic tube}
\ex{The patient was fed through a tube inserted into his stomach.}

\begin{origin}
17 世纪借自拉丁语 \textit{tubus}（管），可能与日耳曼语族词形有关。同源词 \textbf{tubular}（管状的；管乐器发出的声音：\textit{a tubular bell} 管状铃）、\textbf{tubeless}（无内胎的）。英语 tube 与 tub（木盆；浴缸）实为不同源——tub 来自古英语 \textit{tubbe}，与日耳曼语 *\textit{þub-} 凸起相关。
\end{origin}

\vw{vital}{/ˈvaɪtl/}{adj.}

\sense{adj.}{必不可少的，对……极重要的}
\col{vital to/for sth}
\ex{Water is vital to all forms of life on earth.}
\col{vital role/part}
\ex{She played a vital role in the success of the project.}
\col{of vital importance}
\ex{It is of vital importance that we finish on time.}

\sense{adj.}{维持生命所必需的}
\col{vital organs/functions}
\ex{The heart, lungs and brain are all vital organs.}
\col{vital signs 生命体征}
\ex{The nurse checked the patient's vital signs every hour.}

\begin{lex}{vital / crucial / critical / essential 的"重要"强度谱}
\textbf{vital} 强调"攸关生命/成败"（\textit{vital organs}; \textit{a vital decision}）；\textbf{crucial} 强调"决定性"（\textit{a crucial moment}）；\textbf{critical} 强调"成败所系"（\textit{a critical error}）；\textbf{essential} 强调"根本必要"（\textit{essential reading}）。vital 强度最高，常含"失去即亡"之意。
\end{lex}

\begin{origin}
14 世纪借自拉丁语 \textit{vitalis}（属于生命的；生命的），源于 \textit{vita}（生命）。同源词 \textbf{vitality}（活力）、\textbf{vitalize}（赋予生命）、\textbf{vitamin}（维生素——音译自波兰语"vital amine"生命胺）、\textbf{devitalize}（使失去生命）。"生命力"是 vital 的字面义——重要到失去即死。
\end{origin}

\vw{discipline}{/ˈdɪsəplɪn/}{vt. \& n.}

\sense{n.}{训练，纪律}
\col{military/school discipline}
\ex{The new teacher managed to restore discipline in the classroom within a week.}
\col{keep/maintain discipline}
\ex{It is hard to keep discipline among such a large group of children.}

\sense{n.}{行为准则}
\col{a discipline of ...}
\ex{The monks follow a strict discipline of silence and prayer.}

\sense{n.}{自制力}
\col{the discipline to do sth}
\ex{It takes a lot of discipline to write a novel.}

\sense{n.}{学科}
\col{academic disciplines}
\ex{Psychology and linguistics are both academic disciplines that study human behaviour.}

\sense{vt.}{严格要求（自己）}
\col{discipline oneself to do sth}
\ex{He disciplined himself to get up at five every morning to practise the piano.}

\sense{vt.}{惩罚}
\col{discipline sb for sth}
\ex{The sergeant disciplined the soldiers for failing to salute.}

\begin{gram}{discipline 的多义聚焦："训练"的核心}
\textbf{discipline} 名词涵盖：训练/纪律/准则/自制力/学科；动词涵盖：训练/惩罚。一根主线——"通过规则/惩罚使人形成习惯"：学术学科是知识的训练纪律，军队纪律是行为的训练纪律，惩罚是违反纪律的后果。
\end{gram}

\begin{origin}
13 世纪借自古法语 \textit{descepline}，拉丁语 \textit{disciplina}（教导、训练、纪律），源于 \textit{discipulus}（学生；门徒）。同源词 \textbf{disciple}（信徒；弟子）、\textbf{disciplinary}（纪律的；学科的）、\textbf{self-discipline}（自律）。disciple 与 discipline 形近同源——纪律是"门徒受的训练"。
\end{origin}

\vw{leisure}{/ˈleʒə(r)/}{n.}

\sense{n.}{闲暇，空闲}
\col{at leisure}
\ex{Now that the children are at school, I have more time at leisure.}
\col{at one's leisure 在方便时}
\ex{Please feel free to contact me at your leisure.}
\col{leisure time/activities}
\ex{Reading is one of my favourite leisure activities.}

\begin{gram}{at leisure / at one's leisure：休闲的两种时态}
\textbf{at leisure}（处于闲暇中；不慌不忙：\textit{march at leisure} 从容行军）；\textbf{at one's leisure}（在某人方便时；得空时：\textit{Read it at your leisure}）。前者表状态，后者表时间点。
\end{gram}

\begin{origin}
14 世纪借自古法语 \textit{leisir}（被允许；闲暇），拉丁语 \textit{licere}（被允许）。同源词 \textbf{leisurely}（从容的；悠闲的）、\textbf{license}（许可证；放纵；\textit{licentia} 自由）、\textbf{licentious}（放荡的）。"闲暇"在拉丁词源中是"被允许之事"——中世纪欧洲的休闲是贵族特权。
\end{origin}

\vw{crew}{/kruː/}{n.}

\sense{n.}{全体工作人员；全体乘务人员}
\col{crew member 乘务员}
\ex{The captain and crew members welcomed the passengers aboard.}
\col{a flight/cabin crew 航班/客舱乘务组}
\ex{The flight crew prepared the cabin for landing.}

\sense{n.}{专业团队}
\col{a crew of ...}
\ex{A crew of engineers worked round the clock to repair the damaged bridge.}
\col{ground crew 地勤人员}
\ex{The ground crew prepared the aircraft for its next flight.}

\begin{gram}{crew 的单复数与谓语动词}
\textbf{crew} 作集体名词时谓语动词可用单数或复数：\textit{The crew is/are preparing for take-off.}。现代英语倾向单数谓语（强调"整体"），传统/英式偏好复数谓语（强调"成员"）。同类集合名词：\textit{team, staff, family, committee, audience}。
\end{gram}

\begin{origin}
15 世纪中古英语 \textit{crowe}，中古法语 \textit{creue}（增长、增加），源于古法语 \textit{croistre}（增长），拉丁语 \textit{crescere}（生长）。同源词 \textbf{increase}（增长）、\textbf{decrease}（减少）、\textbf{concrete}（具体的；\textit{con-} 一起 + \textit{crescere} 长）、\textbf{crescent}（月牙；增长之物）、\textbf{recruit}（招募；\textit{re-} 再 + \textit{crescere} $\to$再长出来）。"crew"原义"被加进来的人"——船/机成员的增添。
\end{origin}

\vw{permanent}{/ˈpɜːmənənt/}{adj.}

\sense{adj.}{永久的，永恒的}
\col{a permanent job/residence}
\ex{After years of moving from place to place, they finally found a permanent home.}
\col{permanent address/collection}
\ex{Please state your permanent address on the form.}

\begin{gram}{permanent / perpetual / eternal / everlasting 的"永久"细分}
\textbf{permanent} 强调"长期稳定不改变"（\textit{a permanent job}）；\textbf{perpetual} 强调"持续不断"（\textit{perpetual motion} 永动；常含"无止境"之贬义：\textit{perpetual complaints} 无休止的抱怨）；\textbf{eternal} 强调"无始无终"（\textit{eternal life} 永生；常含形而上学/神学意味）；\textbf{everlasting} 强调"持续极久"（\textit{everlasting glory} 永恒的荣耀；文学性强）。注意 perpetual 比 permanent 更含"持续不断"的动感。
\end{gram}

\begin{origin}
15 世纪借自拉丁语 \textit{permanentem}（持续的、稳定的），是 \textit{permanere}（持续、保留）的现在分词：\textit{per-}（贯穿）+ \textit{manere}（停留、保持）。同源词 \textbf{permanence}（永久性）、\textbf{remain}（剩余；继续）、\textbf{remainder}（剩余物）、\textbf{remanent}（剩余的）、\textbf{immanent}（内在的；\textit{im-} 入 + manere）。"永久"源自"贯穿始终停留"。
\end{origin}

\vw{altogether}{/ˌɔːltəˈɡeðə(r)/}{adv.}

\sense{adv.}{完全}
\ex{That is altogether different from what we expected.}

\sense{adv.}{总共}
\ex{How much will it cost altogether?}
\ex{Altogether, there were fifteen people in the room.}

\sense{adv.}{总之}
\ex{Altogether, the project was a great success.}

\begin{lex}{altogether / all together：书写差异与意义}
\textbf{altogether}（adv. 完全；总共；总之——一个词）；\textbf{all together}（所有人/物一起——两个词）：\textit{They went all together to the cinema.}（他们一起去电影院）vs. \textit{They are altogether mad.}（他们完全疯了）。注意：\textit{all together} 强调"聚集在一起"，\textit{altogether} 强调"完全/总共"。
\end{lex}

\begin{origin}
13 世纪 \textit{al}（全）+ \textit{togeder}（一起），古英语 \textit{eall} + \textit{togædere}。与 \textit{together}（一起）同源，但 altogether 已凝固为一个副词。同源词 \textbf{together}（一起）、\textbf{collect}（收集；\textit{col-} 一起 + \textit{legere} 选）。
\end{origin}

\vw{kit}{/kɪt/}{n.}

\sense{n.}{成套设备，成套工具}
\col{a tool/first-aid/repair kit}
\ex{Every car should carry a basic first-aid kit in the boot.}
\col{kit out (with sth) 装备}
\ex{The soldiers were kitted out with the latest equipment.}

\sense{n.}{配套元件}
\col{a kit of parts}
\ex{The model plane comes as a kit of parts that you assemble yourself.}

\begin{gram}{the whole kit and caboodle：口语"全部"加强式}
\textbf{the whole kit and caboodle}（全部东西；所有人）：\textit{Just bring the whole kit and caboodle on Saturday.}。同类的口语加强短语：\textit{the whole nine yards}（从头到尾；全部；该做的都做了）、\textit{the whole shebang}（整个事情）。
\end{gram}

\begin{origin}
14 世纪中期荷兰语 \textit{kitte}（木桶），与北日耳曼语 \textit{kuttr}（小桶）相关。kit 原义"装工具的木桶"，后引申为"桶中之物"——现代"成套工具"是原始"桶"义的转喻。英语 \textit{kitchen}（厨房）与 kit 形近但不同源。
\end{origin}

\vw{universe}{/ˈjuːnɪvɜːs/}{n.}

\sense{n.}{(the universe) 宇宙，天地万物}
\col{the origin/expansion of the universe}
\ex{The expansion of the universe was first observed by the astronomer Edwin Hubble.}
\col{the creator of the universe}
\ex{Philosophers have long debated whether the universe has a creator.}

\begin{origin}
14 世纪借自古法语 \textit{univers}，拉丁语 \textit{universum}（整体、宇宙），是 \textit{universus}（全部的、整个的）的中性形式：\textit{uni-}（一，单）+ \textit{versus}（转动；\textit{vertere} 转动）。同源词 \textbf{universal}（普遍的；通用的）、\textbf{university}（大学；"宇宙的"引申为"普世知识之所"——\textit{universitas magistrorum et scholarium} 师生共同体）、\textbf{verse}（诗；韵文）、\textbf{version}（版本；译文）、\textbf{convert}（皈依；转换；\textit{con-} 一起 + vertere）、\textbf{reverse}（颠倒）。"宇宙"是"万物转动归一之处"——古罗马人对"统一性"的认知。
\end{origin}

\vw{crucial}{/ˈkruːʃl/}{adj.}

\sense{adj.}{至关重要的，关键性的}
\col{crucial to/for sth}
\ex{This decision is crucial to the future of the company.}
\col{a crucial moment/point/role}
\ex{The Battle of Britain was a crucial moment in the Second World War.}
\col{at a crucial stage}
\ex{The negotiations are now at a crucial stage.}

\begin{lex}{crucial / critical / vital / decisive 的"关键"对比}
见 vital 条。\textbf{crucial} 强调"决定性时刻"（\textit{a crucial experiment} 决定性实验）；\textbf{critical} 强调"成败所系"（\textit{a critical decision}）；\textbf{vital} 强调"攸关生命"（\textit{vital organs}）；\textbf{decisive} 强调"终结犹豫"（\textit{a decisive victory}）。四者皆可译"关键"，但侧重点不同。
\end{lex}

\begin{origin}
18 世纪借自法语 \textit{crucial}，拉丁语 \textit{crucialis}（十字形的），源于 \textit{crux}（十字）。crucial 原是神学术语——"处于十字路口的"，引申为"必须做决定的"。同源词 \textbf{crucifix}（十字架像）、\textbf{crucify}（钉在十字架上）、\textbf{cruise}（巡航；"穿过十字路口"$\to$横越；与 crucial 形近但实源不同——cruise 来自荷兰语 \textit{kruisen} 横切）、\textbf{crusade}（十字军；改革运动）。"关键"在词源上是"站在十字路口"。
\end{origin}

\vw{orbit}{/ˈɔːbɪt/}{vt. \& vi. \& n.}

\sense{vt. \& vi.}{围绕……运动，沿轨道运行}
\col{orbit around/round sth}
\ex{The earth orbits the sun once every 365 days.}
\ex{The satellite will orbit the earth at an altitude of 500 kilometres.}

\sense{n.}{(天体等运行的) 轨道}
\col{in orbit}
\ex{The spacecraft is now in orbit around Mars.}
\col{put sth into orbit}
\ex{Russia was the first country to put a human-made satellite into orbit.}
\col{geostationary orbit 地球同步轨道}
\ex{Many communication satellites are placed in geostationary orbit.}

\begin{origin}
17 世纪借自拉丁语 \textit{orbita}（轨道、圆形轨迹），源于 \textit{orbis}（圆、环、球）。同源词 \textbf{orbital}（轨道的；\textit{orbital road} 环路）、\textbf{orb}（球；天体；轨道；\textit{orbits} 眼睛——\textit{orbis} 圆 $\to$ 眼球圆转）、\textbf{orbiter}（轨道飞行器）。拉丁语 \textit{orbis} 兼指"环"与"球"——环形即球形。
\end{origin}

\vw{merely}{/ˈmɪəli/}{adv.}

\sense{adv.}{仅仅，只不过}
\col{merely ... (not ...) 不只是}
\ex{It is merely a question of time before we find a cure.}
\col{not merely ... but also ... 不仅……而且……}
\ex{She is not merely a talented musician, but also a gifted teacher.}

\begin{gram}{merely / purely / simply / only / just 的"仅仅"语体}
\textbf{merely}（仅仅；偏书面、文学）；\textbf{purely}（纯粹地；强调无杂质）；\textbf{simply}（简单地；强调简易）；\textbf{only}（只有；通用）；\textbf{just}（只是；通用）。merely 比 only 正式，置于句中常表"无足轻重"：\textit{I'm merely a beginner, don't expect too much of me.}。
\end{gram}

\begin{origin}
中古英语 \textit{merely}，由 \textit{mere}（纯粹的；仅）+ \textit{-ly} 派生。mere 14 世纪借自拉丁语 \textit{merus}（纯粹的、未掺杂的）。同源词 \textbf{mere}（仅仅的；\textit{a mere child} 仅仅是个孩子）、\textbf{merit}（功绩；\textit{meritum} 应得之物）、\textbf{merge}（合并）。mere 与 \textit{lake}（湖）同形但实不同源——\textit{lake} 来自拉丁语 \textit{lacus} 湖。
\end{origin}

\vw{administration}{/ədˌmɪnɪˈstreɪʃn/}{n.}

\sense{n.}{行政部门}
\col{the Administration}
\ex{The Administration announced new tax cuts yesterday.}
\col{the Bush/Obama Administration}
\ex{The scandal took place during the Nixon Administration.}

\sense{n.}{管理，行政}
\col{the administration of sth}
\ex{She is responsible for the administration of the company's pension scheme.}
\col{public administration}
\ex{He spent his career in public administration at the local level.}

\begin{origin}
14 世纪借自拉丁语 \textit{administrationem}（管理、执行），源于 \textit{administrare}（管理、服务）：\textit{ad-}（朝向、为）+ \textit{minister}（仆人、伺候者）。同源词 \textbf{administer}（管理；执行；给药）、\textbf{administrative}（行政的；管理的）、\textbf{administrator}（管理者）、\textbf{minister}（部长；大臣；伺候者）、\textbf{ministry}（部；部门）。"管理"在拉丁词源中是"伺候"——大臣本是大仆。
\end{origin}

\vw{agency}{/ˈeɪdʒənsi/}{n.}

\sense{n.}{机构}
\col{a government/UN/UN agency}
\ex{The World Health Organization is a UN agency based in Geneva.}
\col{an advertising/employment agency}
\ex{She found her current job through a recruitment agency.}

\sense{n.}{代理处}
\col{by/through the agency of sb}
\ex{The prisoner was released by the agency of a human-rights lawyer.}

\begin{origin}
17 世纪由 \textit{agent}（代理人）+ \textit{-cy} 抽象名词后缀派生。agent 14 世纪借自拉丁语 \textit{agentem}（做的人；行动者），是 \textit{agere}（做、行动）的现在分词。同源词 \textbf{agent}（代理人；动作者；\textit{a secret agent}）、\textbf{agenda}（议程；\textit{agenda} 是 \textit{agere} 中性复数"待办之事"）、\textbf{act}（行动）、\textbf{actor}（演员）、\textbf{action}（行动）、\textbf{active}（积极的）、\textbf{react}（反应）、\textbf{transact}（办理）、\textbf{exact}（精确的；\textit{ex-} 出 + agere $\to$ 做出$\to$强求出来）。
\end{origin}

\vw{comprise}{/kəmˈpraɪz/}{vt.}

\sense{vt.}{包括，包含}
\col{comprise sth}
\ex{The course comprises four modules and a final examination.}
\ex{The team is comprised of players from over ten different countries.}

\sense{vt.}{组成，构成}
\col{be comprised of}
\ex{The class is comprised mainly of students from non-science backgrounds.}
\col{comprise (no passive) 注意主动语态}
\ex{相比 include 的"包括"义，comprise 严格说不用被动语态——应该说 \textit{The committee comprises five members.}，而非 \textit{Five members are comprised in the committee.}。}

\begin{gram}{comprise / consist of / include / be composed of 的"组成"细分}
\textbf{comprise}（包含；包括——主语是整体，宾语是部分；\textit{The team comprises five players}）；\textbf{consist of}（由……组成——主语是部分，介词 of 后是整体；\textit{The team consists of five players}）；\textbf{include}（包括——主语是部分，宾语是包含项；\textit{The team includes two goalkeepers}）；\textbf{be composed of}（由……组成——主语是整体；\textit{Water is composed of hydrogen and oxygen}）。comprise 严格不用于被动语态，是英语写作常考点。
\end{gram}

\begin{origin}
15 世纪借自法语 \textit{comprendre}（包含、理解），拉丁语 \textit{comprehendere}（抓住、包含、理解）：\textit{com-}（一起）+ \textit{prehendere}（抓住、抓住）。同源词 \textbf{comprehend}（理解；包含）、\textbf{comprehensive}（全面的；综合的）、\textbf{apprehend}（理解；逮捕；\textit{ap-} 朝向 + prehendere $\to$ 抓向$\to$逮捕）、\textbf{prison}（监狱——"被抓住之地"）、\textbf{enterprise}（企业；\textit{enter-} 间 + \textit{prehendere} $\to$在中间抓住的尝试）、\textbf{surprise}（惊讶；\textit{sur-} 上 + \textit{prehendere} $\to$从上面抓过来）。
\end{origin}

\vw{scan}{/skæn/}{vt. \& vi. \& n.}

\sense{vt. \& vi.}{扫描}
\col{scan sth for sth}
\ex{The machine scans the barcode for the product information.}

\sense{vt. \& vi.}{细看}
\col{scan the horizon/sky}
\ex{The sailor scanned the horizon for signs of land.}

\sense{vt. \& vi.}{浏览}
\col{scan (through) sth}
\ex{I scanned through the report, looking for the main conclusions.}

\sense{n.}{扫描检查}
\col{a brain/CAT/ultrasound scan}
\ex{The doctors did a brain scan to find the cause of his headaches.}

\sense{n.}{快速查阅}
\col{a quick scan of sth}
\ex{I had a quick scan of the morning's emails before the meeting.}

\begin{lex}{scan / skim / browse / flip through 的"略读"细分}
\textbf{scan}（扫读；快速寻找特定信息：\textit{scan a list for your name}）；\textbf{skim}（略读；快速抓主旨：\textit{skim the headlines}）；\textbf{browse}（随意翻阅；无明确目的：\textit{browse a bookshop}）；\textbf{flip through}（翻页；最快速率：\textit{flip through a magazine}）。四者皆"非精读"，但目的与节奏不同。
\end{lex}

\begin{origin}
14 世纪借自拉丁语 \textit{scandere}（攀爬、扫描、吟诵诗的格律）。同源词 \textbf{ascend}（上升；\textit{a-} 朝向 + scandere）、\textbf{descend}（下降；\textit{de-} 下 + scandere）、\textbf{scan}（扫描——视觉上"逐点攀爬"）、\textbf{scansion}（诗的格律分析；吟诵需"上上下下"逐音节扫描）、\textbf{transcend}（超越；\textit{tran-} 越过 + scandere $\to$越过之上）。
\end{origin}

\vp{plain to see}

\sense{短语}{显而易见的}
\ex{It is plain to see that the company is in financial trouble.}
\ex{From the expression on her face, it was plain to see that she disagreed.}

\begin{gram}{plain / clearly / obviously / evidently 的"显而易见"语体}
\textbf{plain to see}（显而易见；偏文学/古雅）；\textbf{clearly}（清楚地；通用）；\textbf{obviously}（显然地；含"理应如此"之判断）；\textbf{evidently}（明显地；多指根据证据）。plain to see 是较古雅的书面表达；现代口语更常用 obviously。
\end{gram}

\vw{foundation}{/faʊnˈdeɪʃn/}{n.}

\sense{n.}{基础，根据}
\col{the foundation of/for sth}
\ex{Respect for human rights is the foundation of any democratic society.}
\col{lay the foundation for sth}
\ex{The two leaders' early meetings laid the foundation for a lasting peace.}

\sense{n.}{地基}
\col{the foundations of a building}
\ex{The foundations of the old castle were washed away by the flood.}

\sense{n.}{创办；基金会}
\col{found a foundation 创办基金会}
\ex{The Gates Foundation is one of the largest charitable foundations in the world.}

\begin{gram}{foundation / base / basis / ground 的"基础"细分}
\textbf{foundation}（基础；多用于抽象或建筑地基；\textit{the foundation of a theory} 理论基础）；\textbf{base}（底部；多用于具体底部或军事基地；\textit{military base}）；\textbf{basis}（根据；多用于抽象根据；\textit{on a daily basis}）；\textbf{ground}（根据；理由——多用于正当性论证：\textit{on grounds of...}）。注意 base/basis 是近义词，但 base 多指物理底部，basis 多指抽象根据。
\end{gram}

\begin{origin}
14 世纪借自拉丁语 \textit{fundationem}（打地基、创立），源于 \textit{fundare}（打地基），词根 \textit{fundus}（底部、底）。同源词 \textbf{found}（创立；建立）、\textbf{founder}（创始人；\textit{the founder of the school}）、\textbf{fund}（基金；资金）、\textbf{profound}（深刻的；\textit{pro-} 向前 + fundus $\to$ 向前到底$\to$深及底部的）。"打地基"是"创立"——双义同源。
\end{origin}

\vp{lay the foundation for}

\sense{短语}{为……打下基础}
\col{lay the foundation(s) for sth}
\ex{Newton's work laid the foundation for modern physics.}
\ex{The early pioneers laid the foundations of the settlement.}

\begin{gram}{lay 与 lie 的时态陷阱}
\textbf{lay}（vt. 放置；下蛋）：\textit{lay-laid-laid-laying}；\textbf{lie}（vi. 躺下；位于）：\textit{lie-lay-lain-lying}。两者过去式都是 lay，极易混淆：\textit{She laid the book on the table.}（她把书放在桌上——lay 的过去式）vs. \textit{She lay on the bed for hours.}（她躺在床上几小时——lie 的过去式）。判断：看主语是放东西的"主动者"还是躺着的"承受者"。
\end{gram}

\vw{purse}{/pɜːs/}{n.}

\sense{n.}{资金，财源}
\col{the public purse 国库}
\ex{Tax increases are always unpopular with voters who have to pay more into the public purse.}
\col{hold the purse strings 掌管钱财}
\ex{In most families, it is the mother who holds the purse strings.}

\sense{n.}{钱包}
\col{a leather/plastic purse}
\ex{She put her purse on the counter and took out her credit card.}

\begin{lex}{purse / wallet / handbag 的"钱袋"细分（英美差异显著）}
英式：\textbf{purse}（钱包；装硬币的小袋）、\textbf{wallet}（男士皮夹；装纸币）、\textbf{handbag}（手提包；女士）。美式：\textbf{purse}（女士手提包）、\textbf{wallet}（男/女通用皮夹）。注意英美 purse 含义几乎相反。
\end{lex}

\begin{origin}
古英语 \textit{purs}（小皮袋），晚期拉丁语 \textit{bursa}（皮袋），希腊语 \textit{byrsa}（皮）。同源词 \textbf{disburse}（支付；\textit{dis-} 离开 + \textit{bursa} $\to$ 从钱袋中取出）、\textbf{bursar}（财务主管；大学财务长）、\textbf{reimburse}（偿还；\textit{re-} 回 + im- 入 + bursa $\to$ 装回钱袋）。purse 与 \textit{leather}（皮革）有词源关系——钱包本是皮袋。
\end{origin}

\vw{tale}{/teɪl/}{n.}

\sense{n.}{故事}
\col{a tale of ...}
\ex{The film is a moving tale of love and loss in wartime Paris.}
\col{fairy tales 童话}
\ex{She told the children a fairy tale before bedtime.}
\col{tales of adventure/adventure tales}
\ex{The book is a collection of tales of adventure from around the world.}

\sense{n.}{讲述，叙述}
\col{tell a tale}
\ex{Her face told a tale of long suffering.}

\begin{lex}{tale / story / narrative / account 的"故事"细分}
\textbf{tale}（故事；偏文学/民间传说，常含虚构性：\textit{a fairy tale}）；\textbf{story}（故事；通用；可虚可实：\textit{a true story}）；\textbf{narrative}（叙述；强调结构/手法：\textit{the narrative of the novel}）；\textbf{account}（报道；记述；偏真实：\textit{a first-hand account}）。tale 比 story 文学性强，常带古雅色彩。
\end{lex}

\begin{origin}
古英语 \textit{talu}（数、讲述、故事），与 \textit{tell}（告诉）同源。同源词 \textbf{tell}（告诉；讲述；辨别）、\textbf{tally}（计数；记账；\textit{talu} 计数 $\to$ 符木）、\textbf{talk}（谈话；\textit{talu} + 反复后缀）。"tale" 与 "tell" 是同一词根的分化——"讲出来的"即是"故事"。
\end{origin}

\vw{origin}{/ˈprɪdʒɪn/}{n.}

\sense{n.}{起源，起因}
\col{the origin(s) of sth}
\ex{The origins of the custom are still debated by historians.}
\col{of ... origin 出身于……}
\ex{The wine is of French origin.}
\col{trace sth to its origin}
\ex{The rumour was eventually traced to its origin in a casual remark by the mayor.}

\sense{n.}{出身}
\col{of humble/humble origins}
\ex{Despite her humble origins, she rose to become one of the most powerful figures in the country.}

\begin{origin}
16 世纪借自法语 \textit{origine}，拉丁语 \textit{originem}（起源、出身），是 \textit{oriri}（升起、出现）的派生形式。同源词 \textbf{original}（最初的；原创的）、\textbf{originate}（起源；创立）、\textbf{orient}（东方；\textit{oriens} 升起——太阳升起之处）、\textbf{orientation}（方向；定位）、\textbf{disorient}（迷失方向）、\textbf{aboriginal}（原住民的；\textit{ab-} 从 + \textit{oriri} $\to$ 从起源开始）。origin 字面"升起之处"——太阳/事物升起的地方。
\end{origin}

\vw{shuttle}{/ˈʃʌtl/}{n. \& vi.}

\sense{n.}{航天飞机}
\ex{The space shuttle was launched from Cape Canaveral.}

\sense{n.}{来往于两地之间的航班（或班车、火车）}
\col{a shuttle bus/service/flight}
\ex{There's a shuttle bus from the airport to the city centre every twenty minutes.}

\sense{vi.}{频繁往来（于两地之间）}
\col{shuttle between A and B}
\ex{She has been shuttling between London and Paris for her new job.}

\begin{origin}
古英语 \textit{scytel}（梭子——织布机上左右穿梭的工具），原始日耳曼语 *\textit{skutilaz}。同源词 \textbf{shuttlecock}（羽毛球；毽子）、德语 \textit{Schüssel}（碗——\textit{scut-} 投掷）。"航天飞机"与"航班"之 shuttle 来自织布机的"梭"——两端往复，是其"在两地点之间反复移动"的隐喻。
\end{origin}

\vw{frequency}{/ˈfriːkwənsi/}{n.}

\sense{n.}{发生率，出现率}
\col{with increasing/great frequency}
\ex{Accidents on this stretch of road are happening with increasing frequency.}
\col{the frequency of sth}
\ex{The frequency of earthquakes in this region is about one every ten years.}

\sense{n.}{频繁}
\col{with alarming/disquieting frequency}
\ex{Violent incidents are occurring with alarming frequency in the capital.}

\sense{n.}{频率}
\col{a high/low/ultrasonic frequency}
\ex{Bats use ultrasonic frequencies to navigate and hunt for insects.}

\begin{origin}
16 世纪借自拉丁语 \textit{frequentia}（人群、聚集、频繁），源于 \textit{frequens}（频繁的、拥挤的）。同源词 \textbf{frequent}（频繁的；常去的）、\textbf{frequently}（频繁地）、\textbf{frequent}（vt. 常去；\textit{He frequents that café}）。拉丁词源"人群拥挤"——frequent 形容人多的地方，引申为"经常出现"。
\end{origin}

\vw{hydrogen}{/ˈhaɪdrədʒən/}{n.}

\sense{n.}{氢，氢气}
\col{hydrogen bomb/peroxide}
\ex{Hydrogen bombs are far more powerful than atomic bombs.}
\col{hydrogen fuel cell 氢燃料电池}
\ex{Hydrogen fuel cells are seen as a clean alternative to petrol engines.}

\begin{origin}
18 世纪由希腊语 \textit{ydro-}（水；\textit{hydor} 的词根形式）+ \textit{-gen}（产生；\textit{gen-} 出生）合成，1787 年由拉瓦锡命名。同源词 \textbf{generate}（产生）、\textbf{genus}（属）、\textbf{genetic}（基因的）、\textbf{hydro-}（水的；\textit{hydroelectric} 水电的）、\textbf{hydrant}（消防栓）。"氢"字面"产生水之物"——燃烧时产生水。
\end{origin}

\vw{motive}{/ˈməʊtɪv/}{n.}

\sense{n.}{原因，动机，目的}
\col{a motive for sth}
\ex{The police are still trying to establish a motive for the attack.}
\col{underlying motive 潜在动机}
\ex{Underlying motives are often more revealing than stated ones.}
\col{profit motive 利润动机}
\ex{The profit motive drives much of the world's economic activity.}

\begin{gram}{motive / motivation / motivation 的"动"细分}
\textbf{motive}（动机；具体的、隐藏的、犯罪学/心理学用词：\textit{the motive for the murder}）；\textbf{motivation}（动机；更广义的"驱动力"：\textit{student motivation} 学生学习动力）；\textbf{motive} 作形容词表"发动的"（\textit{motive power} 动力）。注意 motive 是名词+形容词一形两性，motivation 仅作名词。
\end{gram}

\begin{origin}
14 世纪借自古法语 \textit{motif}（促使行动之物），晚期拉丁语 \textit{motivus}（运动的），源于 \textit{motus}（运动），是 \textit{movere}（移动、推动）的过去分词。同源词 \textbf{motive}（动机的）、\textbf{motivate}（激发）、\textbf{motivation}（动机）、\textbf{motion}（运动；提议）、\textbf{motor}（发动机；汽车）、\textbf{motive power}（动力）、\textbf{emotion}（情绪；\textit{e-} 出 + movere $\to$ 移出心外之物）、\textbf{promote}（促进；晋升；\textit{pro-} 前 + movere $\to$ 向前推动）、\textbf{remote}（遥远的；\textit{re-} 回 + movere $\to$ 移开）。
\end{origin}

\vw{investment}{/ɪnˈvestmənt/}{n.}

\sense{n.}{投资}
\col{make an investment}
\ex{They made a substantial investment in property.}
\col{an investment in sth}
\ex{An investment in education is an investment in the future.}
\col{investment income 投资收益}
\ex{She lives off her investment income.}

\sense{n.}{投入}
\col{investment of time/effort/energy}
\ex{The success of the project required a huge investment of time and effort.}

\begin{origin}
16 世纪由 \textit{invest}（投资；授予）派生。invest 14 世纪借自意大利语 \textit{investire}（穿衣、授予、包围），拉丁语 \textit{investire}（穿衣、装饰）：\textit{in-}（入）+ \textit{vestire}（穿衣）。同源词 \textbf{invest}（投资；授予；\textit{invest sb with authority} 授予某人权力）、\textbf{investor}（投资者）、\textbf{vest}（背心；授予；\textit{vested interest} 既得利益）、\textbf{travesty}（歪曲；\textit{tra-} 越过 + \textit{vestire} $\to$ 装扮过头）。"投资"字面"穿衣"——把资金"穿"在某事上，赋予其"权力"。
\end{origin}

\vw{calculate}{/ˈkælkjuleɪt/}{vt.}

\sense{vt.}{计算，核算}
\col{calculate the cost/distance/speed}
\ex{Scientists calculated the speed of light with remarkable accuracy.}
\col{calculate that ...}
\ex{Experts calculate that the project will take at least three years to complete.}

\sense{vt.}{预测}
\col{calculate to do sth}
\ex{The speech was carefully calculated to win over the undecided voters.}

\begin{gram}{calculate / compute / reckon / estimate 的"计算"细分}
\textbf{calculate}（精确计算；用公式：\textit{calculate the orbit}）；\textbf{compute}（计算机式计算；现代：\textit{compute the data}）；\textbf{reckon}（心算；估算：\textit{reckon the cost}）；\textbf{estimate}（粗略估计：\textit{estimate the damage}）。calculate 偏精确，estimate 偏粗略。
\end{gram}

\begin{origin}
16 世纪借自拉丁语 \textit{calculatus}（计算的），是 \textit{calculare}（用小石子计算）的过去分词，源于 \textit{calculus}（小石子；计算）。同源词 \textbf{calculation}（计算）、\textbf{calculator}（计算器）、\textbf{calculus}（微积分；结石——"小石子"的医学转义）、\textbf{recalcitrant}（顽抗的；\textit{re-} 回 + \textit{calx} 脚跟$\to$用脚后跟踢回的）。"计算"源于"用小石子数"——古罗马人用算盘珠算账。
\end{origin}

\vw{coverage}{/ˈkʌvərɪdʒ/}{n.}

\sense{n.}{新闻报道}
\col{media/news/press coverage}
\ex{The story received extensive media coverage for several weeks.}
\col{TV coverage of the event}
\ex{The TV coverage of the event was watched by millions.}

\sense{n.}{覆盖范围；信息范围}
\col{the coverage of a plan/policy/insurance}
\ex{Please read the small print to understand the exact coverage of your insurance policy.}
\col{geographical/wifi coverage}
\ex{The geographical coverage of the network now includes most of Western Europe.}

\begin{gram}{coverage / cover 的名词分叉}
\textbf{cover}（vt. 覆盖；n. 覆盖物；封面；掩护；保险）；\textbf{coverage}（n. 覆盖范围；新闻报道）。同根但 coverage 专指"覆盖之范围或报道之内容"，比 cover 更抽象、专业。常见搭配：\textit{TV coverage}（电视报道）、\textit{health coverage}（医保覆盖）、\textit{full coverage}（全额保险）。
\end{gram}

\begin{origin}
15 世纪由 \textit{cover}（覆盖）+ \textit{-age} 抽象名词后缀派生。cover 13 世纪借自古法语 \textit{covrir}，晚期拉丁语 \textit{coperire}（完全覆盖）：\textit{co-}（完全）+ \textit{operire}（盖上）。同源词 \textbf{cover}（覆盖；掩护；封面；\textit{under cover of darkness} 在夜色掩护下）、\textbf{discover}（发现；\textit{dis-} 离 + \textit{coperire} $\to$ 揭开盖子）、\textbf{uncover}（揭开；揭露）、\textbf{recovery}（恢复）。
\end{origin}

\vw{incredibly}{/ɪnˈkredəbli/}{adv.}

\sense{adv.}{极其，极端地}
\ex{She's an incredibly talented musician.}
\ex{The view from the top was incredibly beautiful.}

\sense{adv.}{令人难以置信地}
\ex{Incredibly, no one was hurt in the accident.}

\begin{gram}{incredibly / unbelievably / tremendously / extremely 的"极度"语体}
\textbf{incredibly}（极其；偏口语化，兼表"难以置信"义）；\textbf{unbelievably}（难以置信地；语气更重）；\textbf{tremendously}（极大地；偏正面：\textit{tremendously popular}）；\textbf{extremely}（极其；通用；中性）。incredibly 在英式口语中已普遍等同于"非常"，不必字面理解"难以置信"。
\end{gram}

\begin{origin}
14 世纪由 \textit{incredible}（难以置信的）+ \textit{-ly} 派生。incredible 借自拉丁语 \textit{incredibilis}（不可相信的）：\textit{in-}（不）+ \textit{credibilis}（可信的），源于 \textit{credere}（相信）。同源词 \textbf{incredible}（难以置信的；极好的——口语褒义化）、\textbf{credible}（可信的）、\textbf{credit}（信用；学分；\textit{on credit} 赊账）、\textbf{creed}（信条）、\textbf{credential}（凭证；\textit{credentia} 信物）。
\end{origin}

\vw{dust}{/dʌst/}{n.}

\sense{n.}{沙土，尘土}
\col{a cloud of dust}
\ex{The car raced past in a cloud of dust.}
\col{dust storm 沙尘暴}
\ex{A massive dust storm swept across the city, blanketing everything in a thin layer of sand.}

\sense{n.}{灰尘，尘埃}
\col{dust particles}
\ex{Sunlight streaming through the window revealed thousands of dust particles in the air.}
\col{gather dust 积灰；被遗忘}
\ex{My old typewriter has been gathering dust in the attic for years.}

\begin{gram}{dust 的动词"掸灰"用法}
\textbf{dust sth}（掸去灰尘：\textit{dust the shelves}）、\textbf{dust off}（掸去……上的灰尘；重新使用：\textit{dust off one's French} 重拾法语）、\textbf{bit the dust}（阵亡；失败；俚语）、\textbf{leave sb in the dust}（远远超过某人）。
\end{gram}

\begin{origin}
古英语 \textit{dust}（尘土），与古高地德语 \textit{tunst}（尘）同源。同源词 \textbf{dusty}（满是灰尘的）、\textbf{duster}（掸子；防尘外衣）、德语 \textit{Dunst}（蒸汽、雾）。dust 与拉丁语 \textit{pulvis}（尘；\textit{powder} 粉；\textit{pulverize} 粉碎）不同源。
\end{origin}

\vw{shelter}{/ˈʃeltə(r)/}{n. \& vt. \& vi.}

\sense{n.}{居所，住处}
\col{seek/take shelter}
\ex{We took shelter from the rain under a large oak tree.}
\col{a shelter from ...}
\ex{The hut provided a welcome shelter from the biting wind.}

\sense{n.}{庇护}
\col{under the shelter of sth}
\ex{The troops advanced under the shelter of a smokescreen.}
\col{shelter from ...}
\ex{The trees gave shelter from the harsh afternoon sun.}

\sense{vt.}{保护，掩蔽}
\col{shelter sb/sth from sth}
\ex{The barn sheltered the family from the worst of the storm.}

\sense{vi.}{躲避}
\col{shelter from sth}
\ex{We sheltered from the rain under a bridge.}

\begin{origin}
16 世纪可能借自 \textit{sheld}（盾；古英语 \textit{scyld} 盾牌）+ \textit{-ter} 后缀，或直接源自日耳曼语 *\textit{skeld-}（盾）。同源词 \textbf{shelter} 的核心隐喻是"盾"——shield 是"盾"本身，shelter 是"盾能提供的保护"。
\end{origin}

\vw{constant}{/ˈkɒnstənt/}{adj.}

\sense{adj.}{固定的，不变的}
\col{a constant speed/temperature}
\ex{The machine operates at a constant speed of 3,000 revolutions per minute.}
\col{a constant value 恒定值}
\ex{Some physical constants, like the speed of light, never change.}

\sense{adj.}{连续发生的，重复的}
\col{constant interruptions/complaints}
\ex{She was distracted by the constant noise from the street.}
\col{a constant stream of ...}
\ex{A constant stream of visitors came to see the exhibition.}

\begin{lex}{constant / continual / continuous / perpetual 的"持续"细分}
\textbf{constant}（恒定的；持续不变：\textit{a constant temperature}）；\textbf{continual}（反复的；中间有间断：\textit{continual interruptions}）；\textbf{continuous}（连续的；中间无间断：\textit{continuous noise}）；\textbf{perpetual}（永久的；常含"恼人的无止境"：\textit{perpetual nagging}）。考试常考点：continual 与 continuous 区别——continual 间歇性，continuous 不间歇。
\end{lex}

\begin{origin}
14 世纪借自拉丁语 \textit{constantem}（坚定的、稳定的），是 \textit{constare}（站在一起、坚持）的现在分词：\textit{con-}（一起）+ \textit{stare}（站立）。同源词 \textbf{constancy}（恒久不变）、\textbf{distant}（远的；\textit{dis-} 离 + stare $\to$ 站得远）、\textbf{instant}（立即的；\textit{in-} 内 + stare $\to$ 站在眼前$\to$紧迫的）、\textbf{substance}（物质；\textit{sub-} 在下 + stare $\to$ 站在下面的基础）、\textbf{obstacle}（障碍；\textit{ob-} 对 + stare $\to$ 站在对面）。
\end{origin}

% ---------------- 选收的未标 * 词条 ----------------

\vw{astronaut}{/ˈæstrənɔːt/}{n.}

\sense{n.}{宇航员，航天员}
\col{astronaut training}
\ex{Astronaut training is extremely demanding, both physically and mentally.}

\begin{origin}
20 世纪由希腊语 \textit{astro-}（星；\textit{astron} 星辰）+ \textit{nautes}（水手；航行者）合成。同源词 \textbf{nautical}（航海的；\textit{nautical mile} 海里）、\textbf{astronavigation}（天文导航；星际航行）、\textbf{argonaut}（阿尔戈英雄；\textit{Argo} 船名）。"宇航员"字面"星际水手"——人类早期宇宙航行被视为"海上航行"的延伸。
\end{origin}

\vw{mechanic}{/məˈkænɪk/}{n.}

\sense{n.}{机械师，技工}
\col{a car/garage mechanic}
\ex{We had to call a mechanic when the car broke down on the motorway.}

\begin{origin}
16 世纪借自拉丁语 \textit{mechanicus}（机械的、手工的），希腊语 \textit{mechanikos}（关于工具的），源于 \textit{mechane}（机器、工具）。同源词 \textbf{machine}（机器）、\textbf{mechanical}（机械的；呆板的）、\textbf{mechanism}（机制；机械装置）、\textbf{mechanize}（使机械化）。希腊语 mechane 最初指剧场吊景的"机关"——戏剧里"神的降临"由机械装置完成。
\end{origin}

\vw{assistant}{/əˈsɪstənt/}{n. \& adj.}

\sense{n.}{助理}
\col{a research/personal/shop assistant}
\ex{She works as a research assistant in the chemistry department.}

\sense{adj.}{助理的，副的}
\col{an assistant manager/professor}
\ex{He was recently promoted to assistant manager.}

\begin{origin}
15 世纪由 \textit{assist}（帮助）+ \textit{-ant} 表人派生。assist 15 世纪借自拉丁语 \textit{assistere}（站在旁边）：\textit{ad-}（朝向）+ \textit{sistere}（站立）。同源词 \textbf{assist}（帮助）、\textbf{assistance}（帮助）、\textbf{consist}（组成；\textit{con-} 一起 + sistere $\to$站在一起$\to$由……组成）、\textbf{insist}（坚持；\textit{in-} 在上 + sistere $\to$站在上面$\to$不动）、\textbf{resist}（抵抗；\textit{re-} 反 + sistere）、\textbf{persist}（坚持；\textit{per-} 贯穿 + sistere）。
\end{origin}

\vw{astronomer}{/əˈstrɒnəmə(r)/}{n.}

\sense{n.}{天文学家}
\col{an amateur astronomer}
\ex{Amateur astronomers have made many important discoveries in modern times.}

\begin{origin}
14 世纪由 \textit{astronomy}（天文学）+ \textit{-er} 派生。astronomy 借自拉丁语 \textit{astronomia}，希腊语 \textit{astronomia}（星辰的排列、法则）：\textit{astro-}（星辰）+ \textit{nomos}（法则）。同源词 \textbf{astronomy}（天文学）、\textbf{astronomical}（天文学的；巨大的——\textit{an astronomical sum}）、\textbf{astrology}（占星术；\textit{-logia} 学问——和 astronomy 同根异支）。
\end{origin}

\vw{radiation}{/ˌreɪdiˈeɪʃn/}{n.}

\sense{n.}{辐射，放射线}
\col{radiation therapy/sickness}
\ex{Radiation therapy is one of the main treatments for certain types of cancer.}
\col{background radiation}
\ex{Background radiation from natural sources is present everywhere on earth.}

\begin{origin}
16 世纪由 \textit{radiate}（辐射；发射）派生。radiate 16 世纪借自拉丁语 \textit{radiare}（发射光线），源于 \textit{radius}（光线、辐条）。同源词 \textbf{radiate}（辐射）、\textbf{radiator}（暖气片；散热器）、\textbf{radial}（辐射状的）、\textbf{ray}（光线）、\textbf{radio}（无线电）。"辐射"字面"沿半径方向射出"——与"光线"同源。
\end{origin}

\vw{pill}{/pɪl/}{n.}

\sense{n.}{药丸，药片}
\col{take a pill}
\ex{She takes a pill every morning for her blood pressure.}
\col{a sleeping/the pill 避孕药}
\ex{The pill is one of the most popular forms of contraception.}
\col{a bitter pill to swallow 不得不吞的苦药}
\ex{Losing the match was a bitter pill to swallow.}

\begin{origin}
15 世纪中古拉丁语 \textit{pila}（球、药丸），拉丁语 \textit{pila}（球）。同源词 \textbf{pillar}（柱；\textit{pila} 柱$\to$细长物）、\textbf{compile}（编辑；\textit{com-} 一起 + pilare 压榨）、\textbf{pilfer}（偷窃；\textit{pilfer} 来自古法语 \textit{pelfrer}——同根异形）。pill 与 pile（堆）实为不同源——pile 来自拉丁语 \textit{pila} 柱基。
\end{origin}

\vw{telescope}{/ˈtelɪskəʊp/}{n.}

\sense{n.}{望远镜}
\col{a radio/optical/space telescope}
\ex{The James Webb Space Telescope is the largest optical telescope in space.}

\begin{origin}
17 世纪由希腊语 \textit{tele-}（远）+ \textit{skopein}（看）合成。同源词 \textbf{microscope}（显微镜；\textit{micro-} 小）、\textbf{periscope}（潜望镜；\textit{peri-} 周围）、\textbf{stethoscope}（听诊器；\textit{stethos} 胸）、\textbf{scope}（范围；视野）。telescope 字面"远观之物"。
\end{origin}

\vw{rocket}{/ˈrɒkɪt/}{n.}

\sense{n.}{火箭}
\col{launch a rocket}
\ex{The country successfully launched a rocket into orbit.}
\col{rocket science 火箭学（口语：难事）}
\ex{It doesn't take rocket science to figure out what's wrong.}

\begin{origin}
17 世纪借自意大利语 \textit{rocchetta}（小纺锤），\textit{rocca}（纺锤、岩石）的指小词。同源词 \textbf{rock}（岩石；摇动；与 \textit{rocca} 纺锤不同源——rock 来自古英语 \textit{rocc} 岩石）、法语 \textit{rochet}（主教披肩）。火箭的"纺锤"义源自其形状——细长如纺锤。
\end{origin}

\vw{nuclear}{/ˈnjuːkliə(r)/}{adj.}

\sense{adj.}{核能的，原子能的；核武器的}
\col{nuclear energy/weapons}
\ex{There is widespread concern about the spread of nuclear weapons.}
\col{nuclear family 核心家庭}
\ex{The nuclear family consists of two parents and their children.}

\begin{origin}
19 世纪借自法语 \textit{nucléaire}，拉丁语 \textit{nucleus}（核、仁）。同源词 \textbf{nucleus}（核；核心）、\textbf{nucleic acid}（核酸）、\textbf{nucleate}（成核）、\textbf{nuclear physics}（核物理）。nucleus 拉丁本义"坚果的核"——"核桃仁"。19 世纪科学家借以命名细胞核与原子核。
\end{origin}

\end{document}
